\section{Related work}
\label{sec:related}

The neighbouring literature is strongest exactly where this paper does not compete. Every mechanism discussed below is treated as a component the design may absorb, not as a rival: a stronger proposer, attributor or acceptor makes the authority boundary drawn here more useful, not less.

Archives and self-modification came first. ADAS automates agent-system design through a meta-agent and an archive \cite{adas2024}; Darwin G\"odel Machine performs open-ended self-code modification with empirical validation, an archive and sandboxing \cite{dgm2025}; and ADIAS makes persistent issue identity, lifecycle, evidence and intervention-outcome history a first-class optimization state \cite{adias2026}. Persistent issue state, an archive of candidates and sandboxed validation are therefore inherited here rather than introduced.

Failure attribution is a second neighbourhood, and a crowded one. SAGE uses multi-hypothesis failure attribution to generate evidence-grounded explanations and routes verified root causes to intervention levels \cite{sagemhfa2026}. CausalFlow uses counterfactual intervention to estimate causal responsibility and to generate minimal repairs \cite{causalflow2026}. Failure-driven self-improvement converts failed computer-use trajectories into diagnosed code-level improvements \cite{failurelearning2026}. Each of these answers \emph{which cause explains the failure}. The question they leave open is what an attributed cause is then permitted to authorize, and on whose evidence that authorization rests.

Acceptance-specific work removes another possible novelty route. PACE casts repeated incumbent-versus-candidate commits under adaptive self-evolution as anytime-valid paired sequential testing, so optional-stopping-safe commit acceptance is not a contribution of this paper \cite{paceacceptance2026}. SEA independently confines modification around a frozen base and admits changes through anytime-valid certificates \cite{sea2026}. This paper therefore treats acceptance statistics and verifier gates as absorbed mechanism pressure, not as its residual contribution.

The broader recursive-self-improvement literature raises the bar once more. A 2026 survey of 1,250 recent papers separates bounded self-refinement from stronger recursive self-improvement and places self-evaluation on a verification hierarchy, with demonstrated improvement strength tracking stronger external or formal verification \cite{rsisurvey2026}. PAST-Bench independently asks whether retained experience actually causes later-task improvement under matched experience-on/off conditions and whether the gain follows the intended save--retrieve--update pathway \cite{pastbench2026}. SEVA reports that repeated self-evolution can produce benchmark specialists that improve on one benchmark while regressing on another \cite{seva2026}. These results make generic claims such as ``the system learns from failure'', ``the system improves itself'', ``the candidate passed a verifier'', or ``experience caused later gains'' mature targets rather than a defensible a breakthrough here.

The surviving this paper residual is a different scientific question: \emph{who has authority to turn an internally generated diagnosis or proposal into an adopted scientific method?} Self-ORION separates proposal authority from adoption authority. Its internal controller may diagnose, nominate a revision, select a computation, recommend containment or compose a next-step report, but the registered control object grants none of those outputs revision, adoption, promotion, merge or global-stop authority. The three retained hidden-cause errors therefore become a non-vacuous test condition for the architecture: fallible internal diagnosis can guide investigation without becoming a self-issued adoption certificate. This bounded non-self-promotion result is the paper's present contribution. Transferable fresh-task improvement against strong self-evolving baselines remains a separate prospective extension and is not required to establish the authority-separation theorem.
Self-evolving coding systems press on the same boundary from the engineering side. MOSS rewrites agent source from production-failure evidence, validates candidate images by replay, and places deployment behind user consent and rollback checks \cite{moss2026}. Ratchet manages a self-authored skill library with outcome-driven retirement and bounded capacity while measuring held-out transfer \cite{ratchet2026}, and Verifier-as-Gatekeeper shows that harmful skill admission contaminates later evolution, testing heterogeneous pre-commit gating together with frozen-pool transfer \cite{vag2026}. MEGA accumulates durable optimization assets in a typed wisdom graph and uses controlled evaluation to attribute improvement effects across optimization cycles \cite{mega2026}, while RewardHackingAgents makes evaluator tampering and held-out leakage measurable through patch and access telemetry and evaluator locking \cite{rewardhackingagents2026}. A current survey of self-evolving coding agents treats feedback reliability, benchmark overfitting, safety, maintainability, cost and generalization as field-level concerns \cite{selfevolvingcoding2026}; the position taken here is that those concerns are to be measured, not cited as motivation.

Acceptance statistics form a third neighbourhood. PACE casts repeated incumbent-versus-candidate commits under adaptive self-evolution as anytime-valid paired sequential testing \cite{paceacceptance2026}, and SEA admits self-modifications through anytime-valid certificates around a frozen base \cite{sea2026}. An anytime-valid acceptor answers when an observed difference is real under optional stopping. It does not answer which evidence the difference is allowed to be computed over, and it is adopted here as a within-stage device rather than as the acceptance rule itself.

Finally, transfer measurement supplies the outcome discipline. PAST-Bench evaluates whether retained experience causes later-task improvement under matched experience-on/off conditions \cite{pastbench2026}, which motivates pathway-sensitive evaluation rather than treating memory presence as learning. SEVA reports that repeated self-evolution can produce benchmark specialists that improve on one benchmark while regressing on another \cite{seva2026}, which is why replay evidence and fresh-transfer evidence are measured and reported separately below rather than averaged.

What remains after all of this is narrower and explicitly compositional. A persistent development issue may license a method change only after evidence-bound causal discrimination and invention-readiness checks; acceptance additionally requires non-compensatory staged evidence in the order static checks, replay, independent protected fresh transfer and non-harm, and protected assurance, with immutable retention of negative, null and harmful history and of compromise telemetry, with evaluator and holdout custody outside challenger write authority, and with self-certification and merge authority structurally absent. Promotion remains a host decision. Whether that composition yields more transferable and less compromise-prone improvement than simpler self-edit or acceptance systems is an empirical question; architecture alone cannot settle it, and the evidence reported here does not.
